\documentclass{article}
\usepackage{PRIMEarxiv} % Core package for the PrimeAI Template
\usepackage{amsmath, amssymb, amsthm, bm, dcolumn} % Add amsthm here for the proof environment
\usepackage[numbers,sort&compress]{natbib} % Natbib for citations
\usepackage{graphicx} % For high-quality images
\usepackage{multicol} % Multi-column figures/tables
\usepackage{hyperref} % Hyperlinks
\usepackage{listings} % Code listings
\usepackage{authblk} % For structured affiliations
% Define theorem style (optional)
\theoremstyle{plain}
\newtheorem{theorem}{Theorem}
\renewcommand{\qedsymbol}{}

% Custom commands for primes and qubit encoding
\newcommand{\primeQubit}[1]{|p_{#1}\rangle}
\newcommand{\primeGate}[1]{U_{p_{#1}}}

\usepackage{wrapfig}
\usepackage[pscoord]{eso-pic}
\usepackage[fulladjust]{marginnote}
\reversemarginpar

% Typesetting improvements without footnote patching
\usepackage[protrusion=true, expansion=true, tracking=false]{microtype}
\microtypecontext{spacing=nonfrench}

% Line numbers
\usepackage[right]{lineno}

% Text layout - adjust as needed
\raggedright
\setlength{\parindent}{0.5cm}
\textwidth 5.25in 
\textheight 8.75in

% Set double spacing
\usepackage{setspace} 
\doublespacing

% Adjust width for specific content
\usepackage{changepage}

% Adjust caption style
\usepackage[aboveskip=1pt,labelfont=bf,labelsep=period,singlelinecheck=off]{caption}

% Remove brackets from references
\makeatletter
\renewcommand{\@biblabel}[1]{\quad#1.}
\makeatother

% Header, footer, and page numbers
\usepackage{lastpage,fancyhdr}
\pagestyle{fancy}
\fancyhf{}
\fancyhead[L]{Preprint - PrimeAI Enhanced Template}
\fancyfoot[C]{\scriptsize Multiplicity Theory © 2024 Ryan Van Gelder - Citizen Gardens \\ Licensed Under MIT and CC BY-NC-SA 4.0.}
\fancyfoot[R]{Page \thepage\ of \pageref{LastPage}}
\renewcommand{\footrule}{\hrule height 2pt \vspace{2mm}}

\begin{document}

\title{Modeling Radiation Fields in a Cavity Using Multiplicity Theory}

\author{Ryan O. Van Gelder \\
Citizen Gardens - The Foundation of Multiplicity \\
\texttt{info@citizengardens.org}}
\maketitle
\begin{abstract}
This paper explores the application of Multiplicity Theory to radiation fields within a cavity. Using a dynamic multiplicity framework, we model eigenmode interactions, quantum feedback, and boundary effects. The integration of prime-based encoding and tensor networks refines traditional approaches, enabling a deeper understanding of mode coupling, emergent phenomena, and quantum effects in confined radiation systems.
\end{abstract}

\section{Mathematical Framework}
We use the \textbf{Dynamic Multiplicity Equation}, enhanced for non-linear interactions, stochasticity, and quantum corrections:
\begin{equation}
    \frac{\partial \rho_k}{\partial t} = \alpha_k(t)\rho_k + \beta_k(t)I_k + \gamma_k(t) \sum_j T_{kj}\rho_j 
    + \lambda(t)\left(\Omega_B(\rho) + \Omega_{FS}(\rho)\right) 
    + \eta_k \rho_k^2 + \zeta_k \sum_{m,n} C_{mn} \rho_m \rho_n + \xi_k(t),
\end{equation}
where:
\begin{itemize}
    \item $\rho_k$: Field amplitude of the $k$-th mode,
    \item $\Omega_B(\rho), \Omega_{FS}(\rho)$: Boundary and field source terms,
    \item $\eta_k$: Non-linear feedback,
    \item $\zeta_k$: Quantum entanglement coefficients,
    \item $\xi_k(t)$: Stochastic noise term.
\end{itemize}

\section{Cavity Radiation Dynamics}
\subsection{Eigenmode Expansion}
The cavity field can be expressed via solutions to the Helmholtz equation:
\begin{equation}
    \nabla^2 \psi_n + k_n^2 \psi_n = 0,
\end{equation}
where $k_n$ are the eigenfrequencies. Multiplicity Theory refines this by encoding eigenmodes with prime-based weights:
\begin{equation}
    M(t) = \sum_{i} \lambda_i \mu_i e^{i\theta_i(t)} v_i,
\end{equation}
with $\lambda_i$ as eigenvalues, $\mu_i$ the multiplicity of states, and $v_i$ eigenvectors.

\subsection{Tensor Interactions for Mode Coupling}
Tensor networks capture the non-linear interaction of modes:
\begin{equation}
    T(t) = \sum_{i,j,k} T_{ijk} \psi_i \psi_j \psi_k,
\end{equation}
representing higher-order dependencies and quantum coherence.

\section{Boundary and Quantum Effects}
\subsection{Boundary Effects}
Boundary conditions influence field evolution through:
\begin{equation}
    \Omega_B(\rho) = \int_{\text{boundary}} \rho(\bm{r}, t) \cdot \nabla \psi_n \, dA.
\end{equation}
This incorporates geometric feedback and quantum tunneling.

\subsection{Quantum Entanglement}
Entangled states are represented as:
\begin{equation}
    \zeta_k \sum_{m,n} C_{mn} \rho_m \rho_n,
\end{equation}
where $C_{mn}$ captures the entanglement structure.

\section{Applications}
\subsection{Quantum Optics}
\textbf{Photon Emission Dynamics:}
\begin{equation}
    P(t) = \sum_k \rho_k^2 e^{-\lambda_k t},
\end{equation}
where $\lambda_k$ represents mode decay rates.

\subsection{Material Science}
\textbf{Cavity-Modified Material Properties:}
Using tensor networks, we model the effect of cavity fields on material excitations:
\begin{equation}
    \mathcal{H}_{\text{int}} = \sum_{i,j} T_{ij} \psi_i \psi_j,
\end{equation}
where $\mathcal{H}_{\text{int}}$ is the interaction Hamiltonian.

\subsection{Quantum Information}
\textbf{Entanglement and Coherence:}
The degree of entanglement is quantified by the von Neumann entropy:
\begin{equation}
    S = -\text{Tr}(\rho \log \rho),
\end{equation}
where $\rho$ is the reduced density matrix of the cavity modes.

\section{Conclusion}
Multiplicity Theory provides a robust framework for modeling radiation fields in cavities, bridging classical and quantum perspectives. By integrating prime-based encoding, tensor networks, and recursive feedback, this approach offers insights into mode coupling, quantum entanglement, and emergent radiation dynamics.

\nocite{*}
\bibliographystyle{unsrt}
\bibliography{references}

\end{document}